\section{Proofs for Section~\ref{sec:growth}}\label{app:growth}

\begin{proof}[Proof of Lemma~\ref{lem:graded}]
(a) Consider the side towards $u_i'$; the other is symmetric. The grid interval
between the nodes at distances $t_k<t_{k+1}$ has length at most $\sigma(t_k)$.
For $i\in I_C$ this is $h+\theta t_k$. For $i\in I_Z$ and a length of at least
two, $\sigma(t_k)\ge2$, so $\sigma(t_k)=\lfloor h+\theta t_k\rfloor\le h+\theta t_k$;
intervals of length one are ignored in $w_i$. The node at distance $t_k$ has
adjacent intervals starting at distances $t_{k-1}$ and $t_k$, both at most
$t_k=\abs{v-c_i}$; the intervals adjacent to the center have length at most $h$.

(b) Let $t_1<\dots<t_m=R$ be the nodes on that side, and $t_0=0$. Every step
except the last is unclipped, and an unclipped step from $t$ has length
$\sigma(t)\ge(\hat h+\theta t)/3$: for $i\in I_C$ trivially; for $i\in I_Z$
with $h\ge1$ because $\lfloor y\rfloor\ge y/2$ for $y\ge1$; for $i\in I_Z$ with
$h<1$, either $\theta t<2$ and $\sigma\ge1>(1+\theta t)/3$, or $\theta t\ge2$
and $\lfloor h+\theta t\rfloor\ge\theta t/2\ge(1+\theta t)/3$. Hence
$t_{k+1}+\hat h/\theta\ge(1+\theta/3)(t_k+\hat h/\theta)$ for $k\le m-2$, and
$(1+\theta/3)^{m-1}\le(t_{m-1}+\hat h/\theta)\theta/\hat h<1+\theta R/\hat h$.
So $m-1<\ln(1+\theta R/\hat h)/\ln(1+\theta/3)$, and
$\ln(1+x)\ge\frac{23}{24}x$ for $0\le x=\theta/3\le\frac1{12}$ gives
$m\le\lceil\frac{72}{23\theta}\ln(1+\theta R/\hat h)\rceil$, which implies the
bound.

(c) Each side has length $R\le h$, and $\sigma(0)\ge R$ (for $i\in I_Z$, $R$ is
an integer and $\lfloor h\rfloor\ge R$), so the first step reaches the endpoint.
\end{proof}

\begin{proof}[Proof of the bound $K(\theta,n_P)$ in Lemma~\ref{lem:states}]
The proof of Lemma~\ref{lem:states} bounds every grid after stage $0$ by
$1+2\lceil\frac4\theta\ln(\frac54+4\sqrt{2n_P})\rceil$ nodes. Since
$\lceil x\rceil<x+1$ and $3\le3/(4\theta)$, this is at most
$\theta^{-1}\varphi(n_P)$, where
$\varphi(m)=\frac34+8\ln(\frac54+4\sqrt{2m})$. Now
$\varphi(1)<16.3\le10\lceil\log_23\rceil$, $\varphi(2)<18.6\le10\log_24$, and
for $m\ge2$, $\varphi'(m)\le4/m$ while
$\frac{d}{dm}10\log_2(m+2)\ge7.2/m$. Hence $\varphi(m)\le10\log_2(m+2)$ for
$m\ge2$, and $\varphi(m)\le10\lceil\log_2(m+2)\rceil$ for all $m\ge1$, so
the bound is at most $K(\theta,n_P)$.
\end{proof}

\paragraph{Bit lengths in Theorem~\ref{thm:approx}.}
Let $\Gamma_X$ be the product of the denominators of the box endpoints and
$\Gamma_F$ that of the coefficients, after substituting fixed coordinates;
both have $O(I)$ bits, and $\abs{\varpi_i},\abs E=O(I)$. Fix a trial $\mu$
and let $\alpha=j_{\max}+\max_i\abs{\varpi_i}+\abs E+\mu K_\mu$. Every node of
every grid of the trial has a denominator dividing $\Gamma_X2^\alpha$. Indeed, the initial
center and the original endpoints have denominators dividing $\Gamma_X$, and a
new node is a center coordinate plus or minus
\[
h_{ij}\,\frac{(1+\theta)^k-1}{\theta}
=2^{E-j-\varpi_i}\,\frac{(2^\mu+1)^k-2^{\mu k}}{2^{\mu(k-1)}},\qquad k\le K_\mu,
\]
or an integer step for $i\in I_Z$, or a clipped endpoint of the current box;
centers and box endpoints are nodes of earlier stages of the same trial or
original data. So the bound does not compound from stage to stage; it
depends on the stage limit only through the term $j_{\max}$ in $\alpha$. Nodes
lie in $\bar X$, so their numerators have $O(I+\alpha)$ bits. Every factor
value, correction $L_iw^2/8$, table entry and message has a denominator
dividing $8\Gamma_F\Gamma_X^24^\alpha$ and is a sum of at most $I+n$ terms,
each a coefficient times a product of at most two nodes or of two
differences of nodes (a message is the value of such a sum at a minimizing
assignment). Over this common denominator every number is an integer of
$b=O(I+\alpha)$ bits, the dynamic program uses only integer additions and
comparisons, and a factor evaluation costs $O(b^2)$ bit operations per
monomial. Since $j_{\max}+1\le c(I+q+1)$ and $\lceil\log_2(n_P+2)\rceil\le I+1$, we
have $\mu K_\mu\le10\mu2^\mu(I+1)$ and $b\le c(1+\mu2^\mu)(I+q+1)$. With
$2^{\mu^*}\le6\sqrt{\bar\kappa}$ and $\mu^*\le3(1+\log_2\bar\kappa)$, the
table work of Theorem~\ref{thm:approx} times $b^2$ gives the stated
$f(p,\bar\kappa)(I+q+1)^5$. The certificate stores the grids of one trial and
is checked by the same dynamic programs. Without growth the schedule still
halts, and every number it forms is bounded as above for its trial.

\begin{proof}[Proof of Lemma~\ref{lem:commonmesh}]
We follow the proofs of Lemmas~\ref{lem:inv} and~\ref{lem:states} with
$a_j=Lnh_j^2$ and $\norm\cdot_{\mathbf L}^2$ replaced by $L\norm\cdot^2$.
Write $h=h_j$ and $\rho=\sqrt{n\kappa}$, and note that $L\le\kappa g$ and
$L\theta^2/2\le L/(16\kappa)\le g/16$. Point growth makes $x^*$ the unique
minimizer, so $x^*\in X^{(j)}$ and $\beta_j\le\OPT$ follow as in
Lemma~\ref{lem:inv}(i).

(G1) For $i\in P$, Lemma~\ref{lem:graded}(a) with mesh $h$ and $L_i\le L$ give
$d_i(z_i)\le\frac{L_i}8(h+\theta\abs{z_i-c_i})^2\le\frac L4h^2
+\frac{L\theta^2}4(z_i-c_i)^2$, and
$(z_i-c_i)^2\le2(z_i-x_i^*)^2+2(c_i-x_i^*)^2$; for $i\notin P$, $d_i=0$.
Summing over $i\in P$ and using $n_P\le n$ gives (G1).

(G2) and (G3). We first show by induction on $j$ that
$\norm{c-x^*}^2\le4\kappa nh^2$. For $j=0$ and any $c^{(0)}\in X$,
$\norm{c^{(0)}-x^*}^2\le\sum_is_i^2\le ns^2=nh_0^2$; no other step uses the
choice of $c^{(0)}$. Let
$Y=\norm{y_j-x^*}^2$. Point growth, $\beta_j\le\OPT$, $Q(y_j)=\beta_j$, (G1)
and the bound on $\norm{c-x^*}^2$ give
\[
gY\le F(y_j)-\OPT\le F(y_j)-\beta_j=D(y_j)
\le\frac L4nh^2+\frac g{16}\bigl(Y+4\kappa nh^2\bigr),
\]
so, with $L\le\kappa g$, $Y\le\frac{16}{15}(\frac14+\frac14)\kappa nh^2
=\frac8{15}\kappa nh^2$. The center of stage $j+1$ is $y_j$ and
$h_{j+1}=h/2$, so $\norm{c^{(j+1)}-x^*}^2\le\frac{32}{15}\kappa nh_{j+1}^2
\le4\kappa nh_{j+1}^2$, which completes the induction and proves the first
two bounds of (G2). By (G1),
\[
D(y_j)\le\frac L4nh^2+\frac L{16\kappa}\Bigl(\frac8{15}+4\Bigr)\kappa nh^2
=\frac8{15}Lnh^2\le\frac9{16}Lnh^2 .
\]
Since $U\le F(y_j)$ after step~(ii) and $\OPT\ge\beta_j$, this proves (G3).
For a grid point $z$ with $Q(z)\le U$ let $Z=\norm{z-x^*}^2$; then
\[
gZ\le F(z)-\OPT=Q(z)-\OPT+D(z)\le U-\beta_j+D(z)\le D(y_j)+D(z)
\le\frac{13}{16}Lnh^2+\frac g{16}\bigl(Z+4\kappa nh^2\bigr),
\]
so $Z\le\frac{16}{15}(\frac{13}{16}+\frac4{16})\kappa nh^2
=\frac{17}{15}\kappa nh^2$, the last bound of (G2).

(G4) Let $i\in P$, let $[a,a']$ be a grid interval of coordinate $i$ retained
at stage $j$, let $v\in\{a,a'\}$ satisfy $m_i(v)\le U$, and let the grid
point $z$ attain $m_i(v)$, so $z_i=v$ and $Q(z)\le U$. By (G2),
$\abs{v-x_i^*}\le\sqrt{17/15}\,\rho h$, $\abs{(y_j)_i-x_i^*}\le\rho h$ and
$\abs{c_i-x_i^*}\le2\rho h$; hence
$\abs{v-(y_j)_i}\le(1+\sqrt{17/15})\rho h$ and
$\abs{v-c_i}\le(2+\sqrt{17/15})\rho h$. An integer interval of length one
lies within $\abs{v-(y_j)_i}+1<4.2\rho h+1$ of $(y_j)_i$. Otherwise the
interval has length at most $w_i(v)\le h+\theta\abs{v-c_i}$ by
Lemma~\ref{lem:graded}(a), and $h\le\rho h$ and
$\theta\rho=\theta\sqrt\kappa\sqrt n\le\rho/\sqrt8$ give that each of its
points is within
\[
\bigl(1+\sqrt{17/15}\bigr)\rho h+\rho h+\frac{2+\sqrt{17/15}}{\sqrt8}\rho h
=\bigl(2+\sqrt{17/15}\bigr)\Bigl(1+\frac1{\sqrt8}\Bigr)\rho h<4.15\rho h
\]
of $(y_j)_i$. So every retained interval, and hence their hull
$X_i^{(j+1)}$, lies within $4.2\rho h+[i\in I_Z]$ of $(y_j)_i$.

(G5) At stage $0$, $h_0=s\ge s_i$, so every grid has at most three nodes by
Lemma~\ref{lem:graded}(c), and $3<8\theta^{-1}$; coordinates outside $P$ have
two nodes. At stage $j+1$ the grid of $i\in P$ is centered at
$(y_j)_i\in X_i^{(j+1)}$ (Proposition~\ref{prop:filter}) with mesh $h/2$, and
by (G4) each side has length at most $R=4.2\rho h+[i\in I_Z]$. With $\hat h$
as in Lemma~\ref{lem:graded}(b) for mesh $h/2$,
$\theta R/\hat h\le8.4\theta\rho+\theta\le8.4\sqrt n/\sqrt8+\frac14
<3\sqrt n+\frac14$, so the grid has at most
$1+2\lceil\frac4\theta\ln(\frac54+3\sqrt n)\rceil\le\theta^{-1}\psi(n)$
nodes, with $\psi(m)=\frac34+8\ln(\frac54+3\sqrt m)$, using
$3\le3/(4\theta)$. Now $\psi(1)<12.4\le8\lceil\log_23\rceil$,
$\psi(2)<14.4\le8\log_24$, and for $m\ge2$, $\psi'(m)\le4/m$ while
$\frac{d}{dm}8\log_2(m+2)\ge5.7/m$. Hence
$\psi(m)\le8\lceil\log_2(m+2)\rceil$ for all $m\ge1$.

\emph{CT with a common mesh.} If $P=\emptyset$ it is CT
(Algorithm~\ref{alg:ct}). Otherwise every recorded path is a valid
certificate by Corollary~\ref{cor:filtercert}, and a successful stage gives
$F(\hat x)=U\le\beta_j+\varepsilon$. For termination, take $\mu\ge2$ with
$\theta=2^{-\mu}\le h_{j_{\max}}/s$, so that $\theta s_i\le h_j$ for all
$j\le j_{\max}$ and $i$. Unclipped steps are at least $h_j$ for
$i\in I_C$, and at least $\max\{1,\lfloor h_j\rfloor\}\ge h_j/2$ for
$i\in I_Z$ with $h_j\ge1$; so a side of length at most $s_i\le h_j/\theta$
takes at most $2/\theta$ steps. For $i\in I_Z$ with $h_j<1$ the grid has at
most $s_i+1<1/\theta+1$ nodes. So every grid has at most
$1+4/\theta<8\theta^{-1}\lceil\log_2(n+2)\rceil$ nodes, and the trial is not
aborted. At stage $j_{\max}$, Lemma~\ref{lem:graded}(a) gives
$w_i(v)\le h_{j_{\max}}+\theta s_i\le2h_{j_{\max}}$ for every node $v$, so
$U-\beta_{j_{\max}}\le D(y_{j_{\max}})\le\sum_{i\in P}L_ih_{j_{\max}}^2/2
\le\frac12Lnh_{j_{\max}}^2\le\frac89\varepsilon$, and the trial succeeds.
Under point growth, the trial $\mu=\max\{2,\lceil\log_4(8\kappa)\rceil\}$ has $8\kappa\theta^2\le1$;
by (G5) it is not aborted, and by (G3) it succeeds at stage $j_{\max}$ at the
latest, because
$\frac9{16}Lnh_{j_{\max}}^2=\frac9{16}Lns^24^{-j_{\max}}\le\varepsilon$.
\end{proof}

\begin{proof}[Proof of Proposition~\ref{prop:sharp}]
The Hessian is block diagonal with blocks
$\bigl(\begin{smallmatrix}\Lambda&\Lambda-2g\\\Lambda-2g&\Lambda\end{smallmatrix}\bigr)$,
whose eigenvalues $2\Lambda-2g$ and $2g$ are positive. Hence
$F(x)\ge g\norm x^2$ with equality along $u_k=-v_k$, so $x^*=0$, $\OPT=0$ and
$g$ is the largest growth constant. Each diagonal entry is $\Lambda$. As $g$
decreases from $\Lambda/2$ to $0$, $\kappa=\Lambda/g$ increases continuously
from $2$ to $\infty$.

The corrected objective separates over pairs, $Q=\sum_kQ_k(u_k,v_k)$ with
$Q_k(u,v)=\frac\Lambda2(u^2+v^2)+(\Lambda-2g)uv-\frac\Lambda8(w(u)^2+w(v)^2)$,
where $w$ denotes $w_i$ in the relevant coordinate. Therefore the min-marginal
of $u_1$ is $m_{u_1}(a)=\min_vQ_1(a,v)+\sum_{k\ge2}\min Q_k$, the minima taken
over the grids. By hypothesis $w(0)\ge h$ in every coordinate, so
$\min Q_k\le Q_k(0,0)\le-\Lambda h^2/4$ for each $k\ge2$, and
$\sum_{k\ge2}\min Q_k\le-(m-1)\Lambda h^2/4=-(n-2)\Lambda h^2/8$.

Fix a node $a$ with $|a|\le R_0$ and put $t^*=-(\Lambda-2g)a/\Lambda$; then
$|t^*|\le|a|\le R_0$, so $t^*$ lies in the box of $v_1$ and in some grid
interval of $G_{v_1}$, of length $\lambda$ say. Let $v_a$ be the endpoint of
that interval nearest to $t^*$. Then $|v_a-t^*|\le\lambda/2$ and
$w(v_a)\ge\lambda$. Completing the square,
$\frac\Lambda2v^2+(\Lambda-2g)av=\frac\Lambda2(v-t^*)^2
-\frac{(\Lambda-2g)^2a^2}{2\Lambda}$, so
\[
Q_1(a,v_a)\le\frac{\Lambda^2-(\Lambda-2g)^2}{2\Lambda}a^2
+\frac\Lambda2\Bigl(\frac\lambda2\Bigr)^2-\frac\Lambda8\lambda^2
=\frac{2g(\Lambda-g)}{\Lambda}a^2 ,
\]
where we also dropped the nonpositive term $-\frac\Lambda8w(a)^2$. Thus
\[
m_{u_1}(a)\le\frac{2g(\Lambda-g)}{\Lambda}a^2-\frac{(n-2)\Lambda h^2}8,
\]
which is at most $0\le U$ when
$a^2\le(n-2)\Lambda^2h^2/(16g(\Lambda-g))=(n-2)h^2\kappa^2/(16(\kappa-1))$.
Because $\kappa^2/(\kappa-1)\ge\kappa$, the hypothesis
$a^2\le(n-2)\kappa h^2/16$ suffices. The coordinates $v_k$ and the other
pairs are symmetric. An interval with an endpoint $a$ satisfying
$m_i(a)\le U$ passes the filtering test~\eqref{eq:retain}.
\end{proof}

\begin{proof}[Proof of Corollary~\ref{cor:uniformgrid}]
\emph{Uniform rule, center.} Suppose a stage is centered at $0$ with box
$X'\ni0$. Its grid in coordinate $i$ consists of the multiples of $h_j$ in
$X'_i$ and the endpoints of $X'_i$. All grid intervals have length at most
$h_j$, and $w_i(0)$ is at least the length of every interval: if some
interval is not adjacent to $0$, the side of $0$ that contains it is longer
than $h_j$, so the interval between $0$ and $\pm h_j$ on that side has length
$h_j$. Hence $w_i(v)\le w_i(0)$ for all nodes $v$, so $D(x)\le D(0)$ and, for
every grid point $x\ne0$, $Q(x)-Q(0)\ge F(x)-F(0)\ge g\norm x^2>0$. The
corrected minimizer is $0$, and so is the next center; by induction every
stage is centered at $0$.

\emph{Uniform rule, filtered box.} Put $R_j=\min\{1,2(k_0+1)h_j\}$; we show
by induction that $X^{(j)}\supseteq[-R_j,R_j]^n$. This holds for $j=0$,
since $X^{(0)}=[-1,1]^n$. At stage $0$ ($h_0=2$) the grid is $\{-1,0,1\}$ in
every coordinate, with $w_i(0)=1$; by Proposition~\ref{prop:sharp} with $h=1$
and $R_0=1$ the node $0$ passes, so both intervals are retained and
$X^{(1)}=[-1,1]^n\supseteq[-R_1,R_1]^n$. Let $j\ge1$. Then $h_j\le1$,
hence $h_j\le R_j$, so $\pm h_j$ are nodes and $w_i(0)=h_j$.
Proposition~\ref{prop:sharp} with $h=h_j$ and $R_0=R_j$ retains every
interval with an endpoint $kh_j$, where $|k|\le k_0$ and $|k|h_j\le R_j$. If
$(k_0+1)h_j\le1$, then $(k_0+1)h_j\le R_j$, and the retained intervals
$[kh_j,(k+1)h_j]$, $-k_0-1\le k\le k_0$, give
$X^{(j+1)}\supseteq[-(k_0+1)h_j,(k_0+1)h_j]^n=[-R_{j+1},R_{j+1}]^n$. If
$(k_0+1)h_j>1$, then $R_j=1$ and $X^{(j)}=[-1,1]^n$; moreover $1/h_j$ is an
integer with $1/h_j<k_0+1$, so the nodes $\pm1=\pm(1/h_j)h_j$ pass, the
intervals $[1-h_j,1]$ and $[-1,-1+h_j]$ are retained, and
$X^{(j+1)}=[-1,1]^n\supseteq[-R_{j+1},R_{j+1}]^n$. This completes the
induction. When $(k_0+1)h_j\le1$ and stage $j+1$ is run, its grid is
centered at $0$ with spacing $h_j/2$ and contains every multiple of $h_j/2$ in
$[-(k_0+1)h_j,(k_0+1)h_j]$, that is, $4(k_0+1)+1$ nodes; and
$k_0+1\ge\sqrt{(n-2)\kappa/16}$ gives $4k_0+5\ge\sqrt{(n-2)\kappa}+1$.

\emph{Graded rule.} Apply Proposition~\ref{prop:sharp} with $h=h_j$ and
$R_0=R$. Every point $t\in[-R,R]$ lies in a grid interval $[a,b]$; if
$t\ge0$ then $0\le a\le R$ (the node $0$ lies in the grid), and if $t<0$ then
$-R\le b\le0$, so the interval has an endpoint of absolute value at most $R$
and is retained. Hence the filtered box contains $[-R,R]$ in every
coordinate. The next grid starts from its center $c'$ inside this box and
uses unclipped steps $h'+\theta t$ with $h'=h_j/2$, so after $k$ steps the
distance from $c'$ is at most $h'((1+\theta)^k-1)/\theta$. One side of $c'$
has length at least $R$ within the box and must be covered; hence at least
$\ln(1+\theta R/h')/\ln(1+\theta)$ nodes lie on that side, besides $c'$.
\end{proof}

\begin{proof}[Proof of the claims in Example~\ref{ex:family}]
(a) Let $d=(1-u,1-v)\in[0,1]^2$. At $(1,1)$ the gradient of
$u^2+v^2-4uv+\frac{u+v}4$ is $(-\frac74,-\frac74)$ and its quadratic part is
$d_1^2+d_2^2-4d_1d_2\ge-(d_1^2+d_2^2)$, so, using $d_i\ge d_i^2$,
\[
u^2+v^2-4uv+\tfrac{u+v}4+\tfrac32\ge\tfrac74(d_1+d_2)-(d_1^2+d_2^2)
\ge\tfrac34(d_1^2+d_2^2).
\]
With $(\omega-\frac12)^2\le2(\omega-\frac u2)^2+\frac12d_1^2$ this gives
$\phi-\phi^*\ge\frac34(d_1^2+d_2^2)+(\omega-\frac u2)^2\ge\frac12\bigl(d_1^2+d_2^2
+(\omega-\frac12)^2\bigr)$. The coupling is nonnegative and vanishes at $x^*$;
summing over blocks gives (a). (b) The diagonal Hessian entries are
$\frac52+\frac{2\deg(b)}{16\Delta_{\mathcal G}}\le\frac{21}8$ for $u_b$ and
$2$ for $v_b,\omega_b$; with (a), $\kappa\le\frac{21}8/\frac12$.
(c) At such a point $\bar x$ the coupling gradient in $u_b$ has absolute
value at most $\frac{2\deg(b)}{16\Delta_{\mathcal G}}\le\frac18$. Hence
$\partial_{u_b}F\ge\frac18$ and $\partial_{v_b}F=\frac14$ at $u_b=v_b=0$,
$\partial_{u_b}F\le-\frac74+\frac18$ and $\partial_{v_b}F=-\frac74$ at
$u_b=v_b=1$, and $\partial_{\omega_b}F=0$. For feasible $x$ let $d=x-\bar x$
and $\delta=\sum_b(\abs{d_{u_b}}+\abs{d_{v_b}})$. The first-order term is at
least $\delta/8$; the second-order term is at least
$\sum_b(d_{\omega_b}-\frac12d_{u_b})^2-\delta^2$. Thus
$F(x)-F(\bar x)\ge\delta/8-\delta^2+\sum_b(d_{\omega_b}-\frac12d_{u_b})^2>0$ for
$0<\delta<\frac18$, and $F(x)-F(\bar x)=\sum_bd_{\omega_b}^2>0$ for $\delta=0\ne d$.
(d) Attach to a bag containing $u_b$ a new bag $\{u_b,v_b,\omega_b\}$.
\end{proof}

\begin{proof}[Proof of the claims in Example~\ref{ex:chain}]
With states $\xi_t\in[0,2^t-1]$, controls $z_t\in[0,1]$ and $\xi_0=0$, the
objective of Proposition~\ref{lim:prop:messages} is
$\Psi_m=\xi_m^2+\sum_{t=1}^m(\xi_t-2\xi_{t-1}-z_t)^2
+\frac18\sum_{t=1}^mz_t(1-z_t)$.
Its second derivatives along the coordinates are $2+8=10$ for $\xi_t$ with
$t<m$, $2+2=4$ for $\xi_m$ and $2-\frac14=\frac74$ for $z_t$. Put
$\varsigma_t=2^{-t}\xi_t\in[0,1-2^{-t}]$. The rescaling multiplies the
curvature of the $t$th state by $4^t$, so the rescaled curvatures are
$10\cdot4^t$ for $\varsigma_t$ with $t<m$, $4^{m+1}$ for $\varsigma_m$ and
$\frac74$ for $z_t$,
and $L=4^{m+1}$. Proposition~\ref{lim:prop:messages} shows that
$\Psi_m\ge\frac18(\norm\xi^2+\norm z^2)\ge\frac18(\norm\varsigma^2+\norm z^2)$,
since $\xi_t^2=4^t\varsigma_t^2$; so $\frac18$ is a growth constant of the
rescaled instance. The feasible point with $\varsigma_1=\delta\in(0,\frac12]$
and all other coordinates $0$ has $\xi_1=2\delta$ and
$\Psi_m=(2\delta)^2+(4\delta)^2=20\delta^2$ (here $m\ge2$), so every growth constant
$g$ is at most $20$ and has $L/g\ge4^{m+1}/20=4^m/5$; the constant $\frac18$
gives $L/g=32\cdot4^m$. The weighted condition number is invariant under the
rescaling, so it stays at most $80$.
\end{proof}
